
% RECOMMENDED %%%%%%%%%%%%%%%%%%%%%%%%%%%%%%%%%%%%%%%%%%%%%%%%%%%
\documentclass[graybox]{svmult}

% choose options for [] as required from the list
% in the Reference Guide

\usepackage{mathptmx}       % selects Times Roman as basic font
\usepackage{helvet}         % selects Helvetica as sans-serif font
\usepackage{courier}        % selects Courier as typewriter font
\usepackage{type1cm}        % activate if the above 3 fonts are not available on your system
\usepackage{makeidx}         % allows index generation
\usepackage{graphicx}        % standard LaTeX graphics tool when including figure files
\usepackage{multicol}        % used for the two-column index
\usepackage[bottom]{footmisc}% places footnotes at page bottom
%\usepackage{caption}
\usepackage{subcaption}
\captionsetup{compatibility=false}
\usepackage{lineno}
\usepackage{savefnmark}
\usepackage{comment}
\usepackage{amsmath}
\usepackage{floatrow}

\usepackage[
    backend=biber,
    style=authoryear-icomp,
    sortlocale=de_DE,
    natbib=true,
    url=true, 
    doi=true,
    eprint=false,
    hyperref=true,
    maxnames=1,
    maxalphanames=1,
    maxcitenames=1,
    language=english
]{biblatex}
\addbibresource{main.bib}

\usepackage[
	colorlinks=true,
	linkcolor=black,
	citecolor=black,
	filecolor=black,
	bookmarks=true,
	bookmarksopen=true,
	bookmarksopenlevel=3,
	plainpages=false,
	pdfpagelabels=true,
	pdftitle={A quenched QCD lattice implementation using Fortran coarrays},
	pdfauthor={Nicolas Lang},
	pdfsubject={Lab report},
	pdfkeywords={QCD, Lattice, Fortran, Coarrays}
]{hyperref}

\DefineBibliographyStrings{english}{andothers={et\ al\adddot}}
%\DeclarePunctuationFormat[textcite]{nametitledelim}{\addspace}
\DeclareFieldFormat{pages}{#1}
\makeatletter
\AtBeginDocument{
  \apptocmd\@floatboxreset
    {\let\blx@imc@ifciteseen\@firstoftwo}
    {}
    {\blx@err@patch{floats}}%
}
\makeatother

\floatsetup[table]{capposition=top}


\makeindex             % used for the subject index
                       % please use the style svind.ist with
                       % your makeindex program

\newif\ifPrintNotes
%\PrintNotesfalse to turn off printing of notes
\PrintNotestrue
\ifPrintNotes
  \newcommand{\Nnote}[1]{\textcolor{green}{[Nico: #1]}}
\else
  \newcommand{\Nnote}[1]{\mbox{}}
\fi




%%%%%%%%%%%%%%%%%%%%%%%%%%%%%%%%%%%%%%%%%%%%%%%%%%%%%%%%%%%%%%%%%%%%%%%%%%%%%%%%%%%%%%%%%

\begin{document}

%\linenumbers

\title*{A quenched QCD lattice implementation using Fortran coarrays\\ }
\subtitle{Lab report}
% Use \titlerunning{Short Title} for an abbreviated version of
% your contribution title if the original one is too long
\author{Nicolas Lang}

% Use \authorrunning{Short Title} for an abbreviated version of
% your contribution title if the original one is too long
\institute{Bergische Universit\"at Wuppertal}
%
% Use the package "url.sty" to avoid
% problems with special characters
% used in your e-mail or web address
%
\maketitle


\abstract{A fortran implementation of quenched QCD using coarrays is presented.}


\section{Introduction} \label{sect:introduction}


\section{QCD}


\subsection{The path integral quantization}

\subsection{Discretization on the lattice}

\subsection{The QCD action}



% TODO: perhaps table with Fortran coarray routines?

\section{Fortran implementation}

This section discusses an implementation of the quenched approximation of QCD on the lattice using coarrays for data parallelism.

\subsection{Fortran coarrays}

Coarrays provide a way to distribute data accross multiple program instances while retaining a concise syntax that abstracts the underlying inter-process or network communication. Formerly an extension of Fortran 95, coarrays are now part of the Fortran 2008 standard. Coarrays are variables with one or several \emph{codimensions} and indices referencing data stored by another program instance, called \emph{image}. \emph{Images} represent copies of the program that run asynchronously and can possess further independent local variables. The coarray entries can be accessed from every image by specifying the respective \emph{coindex} representing the image, where the data is stored. The size of the codimension is determined at runtime from the number of running images. The standard provides for additional methods that enable the user to obtain the number of images and the index of the current image and to introduce a break point to synchronize images. The coarray support in the GNU Fortran compiler is based on the \emph{OpenCoarrays} \cite{opencoarrays} library which in turn is based on the Message Passing Interface (MPI).

\subsection{Design choices}

\subsection{Issues}


\section{Sample measurement}



\pagebreak
\printbibliography


\end{document}
